\section{Independent review: original continuous relation completion}\label{independent-review-original-continuous-relation-completion}

Date: 2026-09-13. Scope: CR1--14 in the parent's \texttt{period\_laplacian/CONTINUOUS\_RELATION\_COMPLETION.md}, including the subsequently written CR12a--c pullback construction. This is a written mathematical review, not a numerical or Lean certificate. Only this new independent review directory is written. No source, owner, public, or remote artifact is edited. No new model examples are needed.

\subsection{Disposition and the actual positivity issue}\label{disposition-and-the-actual-positivity-issue}

Accept CR1--11 and CR13--14 with the exact domains and positive-support maps specified there. Accept the revised unconditional CR12a--c and CR12. The original HCD2 assumption is a uniformly positive finite phase family; it is not, by itself, a proof of that assumption for every one-fold phase. Full-order cancellation proves positivity at a sampled selected root, not positivity of every sampled weight or invertibility of every finite one-fold source Gram. The pullback construction resolves that distinction without adding a hypothesis or choosing a new source norm.

For the actual source one can also prove, without uniform phase positivity, that every finite averaged pullback form is positive definite, and that these average forms decrease to zero. The full graph closure of the original class map is \texttt{H\ x\ E}. Proofs and qualifications follow.

\subsection{1. Continuous measure, density, and domains}\label{continuous-measure-density-and-domains}

Retain the original full-order packet \texttt{h}, tensor degree \texttt{k\textgreater{}=1}, \texttt{g=2\ xi}, \texttt{c=k/2}, \texttt{S(u)=c+iu}, and

\texttt{w\_h(u)=\textbar{}(g/h)(1/2+iu)\textbar{}\^{}2/(2\ pi)}, \texttt{m=w\_h\^{}\{*k\}}, \texttt{dnu=m(u)du}.

HD1 and P54 provide the positive exponential moments for this actual measure; P54 also provides the pointwise exponential bound needed for sampling. No mass normalization occurs: \texttt{nu(R)=mu\_h\^{}k}. The quotient \texttt{g/h} is entire and nonzero after full-order removal. Its real-line zeros are isolated, so \texttt{w\_h\textgreater{}0} almost everywhere. For \texttt{k\textgreater{}=2}, integrating the positive-almost-everywhere product in convolution proves \texttt{m(u)\textgreater{}0} for every real \texttt{u}. The inherited boundedness/integrability makes those convolution integrals finite. In every case \texttt{nu} is atomless and positive on every set of positive Lebesgue measure.

For either \texttt{eta=nu} or \texttt{eta=\textbar{}chi(S)\textbar{}\^{}2\ nu}, a smaller positive exponential exponent absorbs the fixed polynomial factor. Every polynomial is in \texttt{L\^{}2(eta)}. With inner product \texttt{integral\ conjugate(f)\ g\ d\ eta}, orthogonality of \texttt{f} to all \texttt{S}-polynomials gives zero moments of \texttt{sigma=conjugate(f)\ eta} in \texttt{u}, since \texttt{u=(S-c)/i} is invertible. Cauchy--Schwarz supplies the exponentially weighted finite variation needed in CR4. Hence \texttt{F(z)=integral\ exp(izu)\ d\ sigma(u)} is holomorphic in the stated strip; dominated differentiation on each smaller strip is valid for every order. All derivatives at zero vanish. The identity theorem gives \texttt{F=0} on the strip.

The Gaussian uniqueness argument in CR4--5 has the correct Fourier signs: \texttt{gamma\_t(x-u)} contributes \texttt{exp(-i\ zeta\ x)\ exp(i\ zeta\ u)}, so its integral against \texttt{sigma} is zero. Finite variation permits Fubini; compactly supported continuous test functions are uniformly approximated by their Gaussian convolutions. Thus \texttt{sigma=0}, hence \texttt{f=0} almost everywhere. This proves both polynomial-density assertions without a moment-problem uniqueness hypothesis being silently added.

Put \texttt{H=L\^{}2(nu)} and \texttt{H\_chi=L\^{}2(\textbar{}chi(S)\textbar{}\^{}2\ nu)}. The map

\texttt{U\_chi:H\_chi\ -\textgreater{}\ H}, \texttt{f\ -\textgreater{}\ chi(S)\ f}

is a linear onto isometry. The inverse \texttt{v/chi(S)} belongs to \texttt{H\_chi} because its weighted squared norm is exactly \texttt{\textbar{}\textbar{}v\textbar{}\textbar{}\_H\^{}2}. The finitely many real zeros of \texttt{chi(S)} have zero \texttt{nu}-mass; inverse values there are irrelevant to both equivalence classes. This is not an assertion that unweighted multiplication by \texttt{chi} is bounded \texttt{H\ -\textgreater{}\ H}. Its weighted domain is essential and is exactly the inherited norm of its relation image, so it does not replace the source metric.

Density in \texttt{H\_chi} and this onto isometry prove \texttt{closure\_H(chi\ C{[}S{]})=H}. Completing the two terms of the inclusion complex \texttt{{[}chi\ C{[}S{]}\ -\textgreater{}\ C{[}S{]}{]}} in their inherited \texttt{H} norms therefore gives \texttt{{[}H\ -\/-I-\/-\textgreater{}\ H{]}}, whose degree-minus-one contraction is \texttt{I}. The original algebraic quotient \texttt{E=C{[}S{]}/(chi)} maps to its zero completed quotient with kernel all of \texttt{E}. In particular completion does not identify the original algebraic quotient with zero before the stated map.

\subsection{2. The same finite canonical forms and the full graph}\label{the-same-finite-canonical-forms-and-the-full-graph}

For \texttt{N\textgreater{}=q-1}, multiplication by the literal nonzero \texttt{chi} is injective on \texttt{P\_\{N-q\}} and has image \texttt{B\_N=ker\ J\_N}. At \texttt{N=q-1} this is the zero relation space. Polynomial evaluation in \texttt{H} is injective because a nonzero polynomial has only finitely many real-line zeros and the density is positive almost everywhere. Thus every finite original \texttt{M\_N} is positive definite.

Let \texttt{s:E\ -\textgreater{}\ P\_\{q-1\}} be the fixed remainder section and let \texttt{Pi\_N} be orthogonal projection in the actual \texttt{H} norm onto \texttt{B\_N}. Then

\texttt{R\_N=(I-Pi\_N)s}, \texttt{J\_N\ R\_N=I\_E}, \texttt{G\_N=R\_N\^{}*\ R\_N}.

The adjoint in this Gram uses the fixed coefficient inner product on \texttt{E} and the actual \texttt{H} inner product. Every lift is \texttt{s\ x+b}, \texttt{b\ in\ B\_N}; orthogonal decomposition gives unique minimality. In coefficient coordinates the unique minimizer is exactly

\texttt{G\_N=(J\_N\ M\_N\^{}\{-1\}J\_N\^{}*)\^{}\{-1\}}, \texttt{R\_N=M\_N\^{}\{-1\}J\_N\^{}*G\_N}.

All inverses here are legitimate: \texttt{M\_N\textgreater{}0} and \texttt{J\_N} is onto. No phase positivity enters these continuous-source formulas.

The \texttt{B\_N} increase and have dense union. Therefore \texttt{Pi\_N\ -\textgreater{}\ I\_H} strongly. Applying this to the finitely many fixed columns of \texttt{s} gives \texttt{R\_N\ -\textgreater{}\ 0} in the operator norm from any fixed finite norm on \texttt{E} to \texttt{H}, and hence \texttt{G\_N\ -\textgreater{}\ 0} in every fixed coordinate matrix norm. The minimum decreases when relations are enlarged, proving \texttt{G\_N\ downarrow\ 0} in Loewner order. Each finite \texttt{G\_N} remains positive definite, and \texttt{det\ G\_N\ -\textgreater{}\ 0} since \texttt{q\textgreater{}=1}. No cutoff or tensor-degree rate follows.

For every \texttt{x\ in\ E}, \texttt{(R\_N\ x,J\ R\_N\ x)\ -\textgreater{}\ (0,x)}. Thus the closure of \texttt{graph\ J} contains \texttt{\{0\}\ x\ E}. It also contains \texttt{(p,Jp)} for every polynomial \texttt{p}, and is a closed linear subspace. Adding \texttt{(0,x-Jp)} shows it contains \texttt{(p,x)} for every polynomial \texttt{p} and every \texttt{x}. Polynomial density now proves the exact stronger assertion

\textbf{\texttt{closure(graph(J:C{[}S{]}\ subset\ H\ -\textgreater{}\ E))\ =\ H\ x\ E}.}

This is full nonclosability, not just an unbounded-evaluation example. It is compatible with \texttt{J\_N\ R\_N=I\_E} at every finite cutoff. Under every fixed invertible coefficient change \texttt{x=W\ x\textquotesingle{}}, the exact maps become \texttt{G\_N\textquotesingle{}=W*G\_NW}, \texttt{R\_N\textquotesingle{}=R\_NW}, \texttt{J\_N\textquotesingle{}=W\^{}\{-1\}J\_N}; all claims persist.

\subsection{3. Unconditional finite phase pullback, including singular phases}\label{unconditional-finite-phase-pullback-including-singular-phases}

Fix the original \texttt{L\textgreater{}0}, put \texttt{u\_n(theta)=(2\ pi\ n+theta)/L}, and retain \texttt{nu\_n(theta)=(2\ pi/L)m(u\_n(theta))}. Let

\texttt{I\_theta=\{n:nu\_n(theta)\textgreater{}0\}}, \texttt{H\_theta=ell\^{}2(I\_theta,nu\_theta)}, \texttt{T\_theta\ P=(P(c+iu\_n(theta)))\_\{n\ in\ I\_theta\}}.

The pointwise exponential estimate P54/CQ5 and its lattice sum CQ6 prove every finite polynomial evaluation is in this original Hilbert space, including when \texttt{k=1}. Zero-weight coordinates are null classes, never inverted evaluation functionals. Define

\texttt{Q\_\{N,theta\}=T\_theta(P\_N)/T\_theta(B\_N)}

with the quotient norm inherited from \texttt{H\_theta}. Both images are finite-dimensional and closed, whether or not \texttt{T\_theta} is injective. The map

\texttt{rho\_\{N,theta\}:E\ -\textgreater{}\ Q\_\{N,theta\}}, \texttt{{[}P{]}\ -\textgreater{}\ {[}T\_theta\ P{]}}

is well defined and onto. Its exact kernel is

\textbf{\texttt{ker\ rho\_\{N,theta\}=J\_N(ker(T\_theta\textbar{}P\_N))}.}

Indeed \texttt{T\_theta\ P=T\_theta\ b}, for some \texttt{b\ in\ B\_N}, exactly when \texttt{P-b} is a sampling-null polynomial with the same class. Conversely such a null representative maps to zero. The unique pullback form is

\texttt{x*G\_\{N,theta\}x=min\_\{b\ in\ B\_N\}\textbar{}\textbar{}T\_theta(sx+b)\textbar{}\textbar{}\^{}2}.

The minimum exists because it is distance to a closed finite-dimensional relation image. Its minimizing observed vector is unique, although a minimizing coefficient polynomial need not be. Its radical is the exact displayed kernel. If \texttt{M\_\{N,theta\}\textgreater{}0}, this is precisely the usual positive canonical Gram \texttt{(J\_N\ M\_\{N,theta\}\^{}\{-1\}J\_N\^{}*)\^{}\{-1\}}. At a singular phase that inverse expression is not asserted.

For each fixed \texttt{x}, the minimum equals the infimum over a countable dense set of complex rational relation coefficients. The norm at each such coefficient is measurable in \texttt{theta}, so the infimum is measurable; polarization makes all matrix entries measurable. Testing the particular original lift \texttt{R\_N\ x} and unfolding the literal phase measure gives

\texttt{integral\ x*G\_\{N,theta\}x\ dtheta/(2\ pi)}

\texttt{\textless{}=\ integral\ \textbar{}\textbar{}T\_theta\ R\_N\ x\textbar{}\textbar{}\^{}2\ dtheta/(2\ pi)}

\texttt{=\ integral\_R\ \textbar{}R\_Nx(c+iu)\textbar{}\^{}2\ m(u)du\ =\ x*G\_Nx}.

The normalization is exact: \texttt{(2\ pi/L)dtheta/(2\ pi)=du} on each lattice interval. The integrals are finite. Off-diagonal entries are integrable by the positive-form inequality \texttt{\textbar{}G\_ij\textbar{}\textless{}=sqrt(G\_ii\ G\_jj)} and Cauchy--Schwarz in the actual phase measure. Thus for every original \texttt{k\textgreater{}=1} and every \texttt{N\textgreater{}=q-1}, without universal phase invertibility,

\textbf{\texttt{0\ \textless{}=\ bar\ G\_N\ :=\ integral\ G\_\{N,theta\}\ dtheta/(2\ pi)\ \textless{}=\ G\_N},}

\textbf{\texttt{D\_N:=G\_N-bar\ G\_N\ \textgreater{}=\ 0}.}

\subsubsection{Exact relation-gap version without coefficient inverse assumptions}\label{exact-relation-gap-version-without-coefficient-inverse-assumptions}

Let \texttt{Pi\_\{N,theta\}} be orthogonal projection in \texttt{H\_theta} onto \texttt{T\_theta(B\_N)}. Since \texttt{T\_theta\ R\_N\ x} and \texttt{T\_theta\ s\ x} differ by that relation image, the minimizing observed lift is \texttt{(I-Pi\_\{N,theta\})T\_theta\ R\_N\ x}. Orthogonal decomposition gives

\texttt{x*D\_Nx\ =\ integral\ \textbar{}\textbar{}Pi\_\{N,theta\}T\_theta\ R\_N\ x\textbar{}\textbar{}\^{}2\ dtheta/(2\ pi)}.

This is an exact positive original-relation norm at every phase; it specializes to HCD13 on the admitted positive family. For an explicit measurable coefficient formula only, let \texttt{B} contain the literal relation columns, \texttt{F\_theta=B*M\_\{N,theta\}B}, and \texttt{d\_theta=B*M\_\{N,theta\}R\_N}. Then the projected relation is \texttt{T\_theta\ B\ F\_theta\^{}+\ d\_theta}, where \texttt{+} is the finite Moore--Penrose pseudoinverse. Here \texttt{d\_theta} lies in the range of \texttt{F\_theta}: a vector in \texttt{ker\ F\_theta} gives the zero observed relation and hence pairs to zero with \texttt{T\_theta\ R\_N}. Thus this expression is precisely the projection, not a modified quotient. Pseudoinverse is Borel measurable, for example as the limit of \texttt{(F\_theta\^{}2+epsilon\ I)\^{}\{-1\}F\_theta}. Its observed squared norm is bounded by that of \texttt{T\_theta\ R\_N}, so the integral exists even if the chosen coefficient representative is not uniformly bounded. At \texttt{B\_N=0} the projection is zero and \texttt{D\_N=0}.

The mean-zero coefficient contraction and coefficientwise averaging of arbitrary fields in HCD5--16 still use the HCD2 uniform-equivalence hypothesis. The unconditional argument above does not silently extend those stronger coefficient-field assertions beyond their stated domain.

\subsubsection{What is unconditional about one-fold positivity}\label{what-is-unconditional-about-one-fold-positivity}

For \texttt{k=1}, the zeros of the actual nonzero entire \texttt{g/h} on the sampled vertical line form a discrete, hence countable, set. Their phases modulo \texttt{2\ pi} form a countable set. Outside it every lattice weight is positive; a nonzero polynomial cannot vanish at all the infinitely many distinct sample points. Therefore \texttt{M\_\{N,theta\}\textgreater{}0} for almost every phase, simultaneously for all finite \texttt{N}. For \texttt{k\textgreater{}=2} the same statement holds at every phase by everywhere-positive convolution.

Consequently, for every nonzero fixed \texttt{x}, \texttt{x*G\_\{N,theta\}x\textgreater{}0} almost everywhere and its integral is positive. This proves the additional actual-source statement

\textbf{\texttt{bar\ G\_N\textgreater{}0} at every finite cutoff, and \texttt{bar\ G\_N\ downarrow\ 0}.}

Monotonicity follows phasewise by enlarging relations; the limit follows from \texttt{bar\ G\_N\textless{}=G\_N}. Also \texttt{0\textless{}=D\_N\textless{}=G\_N} implies \texttt{D\_N-\textgreater{}0}. No monotonicity of \texttt{D\_N}, no zero limit for \texttt{G\_N\^{}\{-1\}D\_N}, and no uniform positive lower bound in \texttt{theta} is inferred. In particular positivity almost everywhere does not supply the HCD2 uniform lower constant.

\subsection{4. Fixed-phase limits, full primary kernels, and surviving masses}\label{fixed-phase-limits-full-primary-kernels-and-surviving-masses}

The exponential sampled moment proof applies on \texttt{I\_theta}, and also after multiplying the measure by \texttt{\textbar{}chi\textbar{}\^{}2}. Thus the sampled polynomial vectors are dense. The image of multiplication by \texttt{chi} is exactly the closed subspace vanishing at

\texttt{Z\_theta=\{n\ in\ I\_theta:chi(c+iu\_n(theta))=0\}}.

Division by \texttt{chi} on the complementary positive atoms proves the reverse inclusion isometrically in the weighted domain. Restriction to \texttt{Z\_theta} is the orthogonal quotient map, with squared norm \texttt{sum\_\{n\ in\ Z\_theta\}\ nu\_n(theta)\textbar{}v\_n\textbar{}\^{}2}. No weight or phase is removed.

Increasing finite relation images have union equal to the evaluated polynomial relation space. Their distances therefore decrease to distance to its closure, including singular finite phase Grams. For every class \texttt{x}, this proves

\texttt{x*G\_\{N,theta\}x\ downarrow\ sum\_\{n\ in\ Z\_theta\}nu\_n(theta)\textbar{}x(S\_n(theta))\textbar{}\^{}2}.

Polarization gives the matrix limit CR13 in the same fixed coefficient frame. Let \texttt{p\_theta} be the squarefree product of the displayed roots. The exact limiting radical is \texttt{(p\_theta)/(chi)}. In a sampled primary block \texttt{C{[}epsilon{]}/(epsilon\^{}r)}, only the constant coefficient is received and the full ideal \texttt{(epsilon)/(epsilon\^{}r)} stays in the kernel. Every unsampled primary block stays entirely in the kernel. This uses the full original \texttt{chi}; it does not replace its multiplicities by a reduced polynomial upstream.

At \texttt{k=1}, the original cyclic minimal polynomial has ideal \texttt{(chi)=(h)}. If \texttt{a\_h} is the original leading coefficient, \texttt{chi=h/a\_h} identifies the monic ideal generator, without replacing the density \texttt{\textbar{}g/h\textbar{}\^{}2/(2\ pi)}. For an actual selected root \texttt{rho} of complete order \texttt{r}, write \texttt{h=(S-rho)\^{}r\ h\_0}, with the unchanged \texttt{h\_0(rho)!=0}. Then

\texttt{(g/h)(rho)=g\^{}\{(r)\}(rho)/(r!\ h\_0(rho))\ !=\ 0},

\texttt{nu\_n(theta)=\textbar{}g\^{}\{(r)\}(rho)\textbar{}\^{}2/(L\ (r!)\^{}2\ \textbar{}h\_0(rho)\textbar{}\^{}2)\textgreater{}0}

whenever \texttt{S\_n(theta)=rho}. Thus no sampled root of the actual one-fold \texttt{chi} loses its receiving line to a zero weight. This does not say that every phase samples such a root or that every other atom has positive weight. It is the exact narrower positivity needed for CR14.

If a displayed sampled root exists, take the original unit class \texttt{x=1} (or any class evaluating to one there). Since \texttt{G\_\{N,theta\}\textgreater{}=K\_theta}, its numerator is at least the retained positive mass. Its continuous denominator is positive at every finite cutoff and tends to zero by CR11. This proves the quotient ratio in CR14 tends to infinity. If \texttt{Z\_theta} is empty there is no such lower bound and no relative rate is asserted. No sampled root is invented.

Only finitely many phases can sample a root of the fixed \texttt{chi}: each root on \texttt{Re\ S=c} specifies its unique phase \texttt{L\ Im\ rho\ mod\ 2\ pi}. The completed quotient field is therefore zero almost everywhere. Its measurable realization is the range in ordinary \texttt{ell\^{}2(Z)} of the diagonal projection with entries \texttt{1\_\{nu\_n(theta)\textgreater{}0\}\ 1\_\{chi(S\_n(theta))=0\}}. Its direct integral is the zero Hilbert space, and the induced map from the whole original \texttt{E} has kernel \texttt{E}. Individual exceptional weighted fibres are not erased pointwise. The retained marked map is \texttt{(label,v)-\textgreater{}(label,0)}, with \texttt{tau-\textgreater{}tau}; represented zero is not external absence.

\subsection{5. Read scope and provenance pins}\label{read-scope-and-provenance-pins}

Read completely: the original and revised root note; the current HD source; the completed-phase density review; the continuous-descent review; and the complete delivered periodized-source research note. Source SHA-256 pins:

\begin{itemize}
\tightlist
\item
  \texttt{source-workspace/work/holonomy\_degree\_cost\_20260913.tex}: \texttt{9A1BA6CD7E8BEE67C6D7868C74F2CF4F9D8C11FA255D9D278C7A2420CF0DACA6}. This current version includes HD10. The parent's earlier \texttt{56c05d...} pin concerns its earlier HD1--9 snapshot, not this current file.
\item
  \texttt{source-workspace/work/holonomy\_completed\_quotient\_density\_review\_20260913.md}: \texttt{20FC066345E4AD5860C0EEB02BB7E446A69CA0826CC8D8F2196B517BA91EA753}.
\item
  \texttt{source-workspace/work/holonomy\_continuous\_descent\_review\_20260913.tex}: \texttt{E283A93B513CB99398CE5F43BBF230D1494BFB55D59243916B1E017852B376C6}.
\item
  \texttt{source-workspace/output/split\_zero\_rh\_tandem\_2026-09-12/sources/web\_periodized\_source\_delivery/Tau\_Periodized\_Source\_Control/RESEARCH\_NOTE.md}: \texttt{F903AE8473EFDFCDE0F45A83878837289D2B0604C1A34903276D7C840FE972B9}.
\end{itemize}

Key source proof locations are HD1--5 and HD6--10; HCD2 for the uniform phase assumption, HCD9--16 for its positive-family gap formulas, HCD17--29 for sampled support, completion and marked maps; CQ4--8 for actual exponential sampled tails, CQ14--22 for the quotient and primary kernel, CQ23--25 for exact one-fold root masses, CQ26--29 for the direct integral; and P8--10, P26a--b, P53--58 in the delivered note.

This review makes no evaluated arithmetic enclosure, uniform source or tensor-degree estimate, public-publication, or proof-assistant claim.
